\documentclass[11pt]{article}
\usepackage{amsmath,amssymb,booktabs}
\usepackage[margin=1in]{geometry}
\begin{document}

\section*{Step 209: Riemann--Siegel \(Z(t)\) Carrier}

Define
\[
  \theta(t)=\Im\log\Gamma(1/4+it/2)-{t\over2}\log\pi,
  \qquad
  Z(t)=e^{i\theta(t)}\zeta(1/2+it).
\]
The carrier is the critical-line \(t\)-axis, modeled as
\(L^2_{\rm loc}(\mathbb R,dt)\) or a weighted \(L^2(\mathbb R,w(t)dt)\)
for windowed probes.  Native probes are point values \(Z(t_n)\), sign changes,
and zero-counting windows.

\subsection*{Real-valuedness is not RH}

For real \(t\), \(\overline{\zeta(1/2+it)}=\zeta(1/2-it)\).  The functional
equation and the definition of \(\theta(t)\) imply
\[
  \overline{Z(t)}=Z(t).
\]
Thus \(Z(t)\) is real-valued on \(\mathbb R\) independently of RH.  The residual
\(\|\Im Z\|^2\) is therefore a degenerate identity and cannot be the RH carrier
residual.

\subsection*{RH-equivalent residual}

Let \(N(T)\) be the Riemann--von Mangoldt count of all nontrivial zeros with
height in \([0,T]\), and let \(N_0(T)\) be the number of real zeros of \(Z(t)\)
in \([0,T]\), counted with multiplicity and endpoint conventions.  Define
\[
  \Delta_Z(T)=N(T)-N_0(T).
\]
Then
\[
  \Delta_Z(T)=0\quad \hbox{for all }T
  \quad\Longleftrightarrow\quad \hbox{RH}.
\]
A Hilbert/distributional version may test \(\Delta_Z\) against a separating
family of Schwartz windows, but this does not remove the target equivalence.

\subsection*{CRCFT audit}

The carrier is classically-RH-equivalent.  It does not expose a CTMT
matrix-element terminus; it states the target zero-count equality directly.
Therefore the mode is CRCFT-TE:
\[
  \boxed{\texttt{V\_Z\_CRCFT\_TE}.}
\]

Classical references: Riemann's Nachlass; Siegel (1932) for the
Riemann--Siegel formula; Edwards (1974, Ch. 7); Odlyzko's numerical
zero computations; the zeta functional equation; and the
Riemann--von Mangoldt zero-count formula.

\end{document}
